\documentclass[12pt,a4paper,oneside]{book}

%%%%% PACKAGES %%%%%

\usepackage{amsmath,amssymb,amsthm,mathtools}
\usepackage{geometry}
\usepackage{microtype}
\usepackage{hyperref}
\usepackage{cleveref}
\usepackage{graphicx}
\usepackage{tikz}
\usepackage{tikz-cd}
\usetikzlibrary{
  arrows.meta,
  positioning,
  calc,
  fit,
  backgrounds,
  decorations.pathmorphing,
  decorations.markings,
  patterns,
  shapes.geometric
}
\usepackage{setspace}
\usepackage{imakeidx}

\geometry{margin=1.2in}
\setstretch{1.15}

\makeindex

\hypersetup{
  colorlinks=true,
  linkcolor=black,
  citecolor=black,
  urlcolor=black
}

%%%%% NUMBERING %%%%%

\numberwithin{equation}{chapter}

%%%%% THEOREM ENVIRONMENTS %%%%%

\theoremstyle{plain}
\newtheorem{theorem}{Theorem}[chapter]
\newtheorem{proposition}[theorem]{Proposition}
\newtheorem{lemma}[theorem]{Lemma}
\newtheorem{corollary}[theorem]{Corollary}

\theoremstyle{definition}
\newtheorem{definition}[theorem]{Definition}
\newtheorem{example}[theorem]{Example}

\theoremstyle{remark}
\newtheorem{remark}[theorem]{Remark}

%%%%% TITLE %%%%%

\title{Constraint Geometry and Semantic Dynamics:\\[6pt]
A Unified Field Theory of Hierarchical Inference,\\
Cognitive Transformation, and Alignment}

\author{Flyxion}

\date{\today}

%%%%% DOCUMENT %%%%%

\begin{document}

\frontmatter

\maketitle

\chapter*{Preface}

The central claim of this monograph is that structure can be
reconstructed from minimal observation, provided that those
observations participate in a coherent constraint network. Across
disciplines as diverse as phylogenetics, machine learning,
neuroscience, and social epistemology, the same principle reappears
in different guises: sparse local measurements accumulate until a
global structure becomes identifiable. What varies across domains is
not the logic of reconstruction but the ontology of the structures
being recovered.

In classical hierarchical clustering the latent structure is a tree.
In information geometry it is a manifold equipped with curvature. In
neuroscience it appears as coordinated transformations across cortical
fields. In artificial intelligence it manifests as the translation of
representational structure into action. Each of these domains can be
interpreted as a particular instantiation of a more general
mathematical architecture in which constraints act as operators
reducing entropy over a space of hypotheses.

This monograph develops that architecture systematically across
fourteen chapters. The argument begins with foundational modeling
assumptions and the state-space formalism that supports all subsequent
constructions. It then develops hierarchical reconstruction from
triplet constraints and its algorithmic implementation. Ultrametric
geometry reveals the metric structure of discrete hierarchies, and a
dedicated chapter on limits of hierarchical structures makes rigorous
the passage from discrete trees to continuous geometric fields via
Gromov--Hausdorff convergence and metric measure space theory.

The geometry of constraint spaces gives the polyhedral structure of
the ultrametric cone explicit form. Information geometry provides the
language of entropy descent and variational inference. The
Relativistic Scalar--Vector Plenum (RSVP) model then introduces the
continuous field-theoretic substrate, whose well-posedness and
dynamical regimes are analyzed before local transformation
operators---the amplitwists---are introduced as the fundamental
mechanisms through which representations are updated.

The algebraic structure of amplitwist cascades is then analyzed
through the lens of Lie theory: cascades form flows on the Lie group
$\mathbb{R}_{>0} \times SO(n)$ and accumulate torsion through
Baker--Campbell--Hausdorff correction terms. The sheaf-theoretic
formulation of constraint propagation follows, revealing that
hierarchical reconstruction and RSVP field inference are both
instances of the search for global sections of a constraint sheaf.
Spectral geometry of constraint operators connects the framework to
graph Laplacians and diffusion processes on hypothesis spaces.
Stochastic extensions of the RSVP dynamics complete the probabilistic
picture via Fokker--Planck equations and free energy minimization.

Worked examples and an empirical observability chapter connect the
abstract formalism to measurable phenomena. The categorical alignment
chapter formalizes the representation--action gap as functorial
non-preservation of structural invariants, with a measure-theoretic
proof that aligned functors form a set of measure zero. The final
synthesis chapter states and proves the Unified Constraint
Reconstruction theorem, which encompasses all preceding results as
special cases.

Throughout the text the emphasis remains on unification. Hierarchical
reconstruction, field dynamics, cortical computation, semantic
evolution, Lie algebra structure, sheaf cohomology, spectral theory,
stochastic analysis, and alignment theory are treated not as
independent problems but as manifestations of the same underlying
mathematical structure: the emergence of order through
constraint-driven entropy reduction.

\tableofcontents

\mainmatter

%%% CHAPTER 1 %%%

\chapter{Foundational Assumptions and Modeling Scope}
\label{ch:foundations}

\section{Motivation}

The purpose of this monograph is to develop a unified mathematical
framework in which complex structures emerge from the accumulation and
propagation of local constraints. The central claim is that
hierarchical inference, geometric reconstruction, field dynamics,
semantic transformation, and alignment theory can all be understood
as instances of a common structural principle: global organization
arises as the intersection of constraint manifolds acting on an
appropriate state space.

\begin{figure}[htbp]
\centering
\begin{tikzpicture}[
  >=Latex, scale=1,
  every node/.style={font=\small},
  statebox/.style={draw, rounded corners, minimum width=2.8cm,
    minimum height=1cm, align=center},
  constraint/.style={draw, circle, minimum size=9mm, align=center},
  arrow/.style={->, thick}
]
\node[statebox] (X) at (0,0) {$\mathcal{X}$\\State space};
\node[constraint] (C1) at (4,1.5) {$C_1$};
\node[constraint] (C2) at (4,0) {$C_2$};
\node[constraint] (C3) at (4,-1.5) {$C_3$};
\node[statebox, minimum width=3.3cm] (I) at (8,0)
  {$\bigcap_i C_i$\\Admissible region};
\draw[arrow] (X) -- (C1);
\draw[arrow] (X) -- (C2);
\draw[arrow] (X) -- (C3);
\draw[arrow] (C1) -- (I);
\draw[arrow] (C2) -- (I);
\draw[arrow] (C3) -- (I);
\node at (2,2.3) {observations};
\node at (8,-1.8) {$H(p)$ decreases as constraints accumulate};
\end{tikzpicture}
\caption{Foundational architecture. A state space $\mathcal{X}$ is
progressively restricted by local constraints $C_i$, producing an
admissible region $\bigcap_i C_i$. Inference is interpreted as
entropy reduction over this constrained state space.}
\label{fig:foundational-constraint-architecture}
\end{figure}

\section{State Spaces and Structural Hypotheses}

Let $\mathcal{X}$ denote the space of possible configurations of a
system. The elements of $\mathcal{X}$ represent candidate structures
consistent with the modeling assumptions.

\begin{definition}[State Space]
A \emph{state space} is a set $\mathcal{X}$ whose elements represent
possible configurations of a system under consideration.
\end{definition}

\section{Observations as Constraints}

\begin{definition}[Constraint]
A \emph{constraint} is a subset $C \subseteq \mathcal{X}$ representing
the set of states consistent with a given observation.
\end{definition}

If a collection of observations $\{o_1, \dots, o_n\}$ is available,
the admissible state space becomes $\mathcal{X}_n = \bigcap_{i=1}^n
C_i$.

\section{Completeness of Constraint Systems}

\begin{definition}[Complete Constraint System]
\label{def:complete-constraint}
A set of constraints $\{C_i\}$ on $\mathcal{X}$ is \emph{complete}
if $\bigcap_i C_i$ contains exactly one element.
\end{definition}

\section{Entropy and Uncertainty}

The uncertainty about the system's configuration is measured by
Shannon entropy $H(p) = -\sum_x p(x)\log p(x)$. Each observation
reduces entropy. Inference is therefore a process of entropy reduction
driven by the accumulation of constraints.

\section{Domains of Application and Scope}

The framework applies to phylogenetics, machine learning, neuroscience,
linguistics, and artificial intelligence. The text focuses on three
components in each domain: the structure of the state space, the
propagation of constraints, and the dynamical mechanisms through which
structures evolve.

\section{Roadmap}

The monograph develops the theory progressively through fourteen
chapters, moving from discrete combinatorial reconstruction through
ultrametric geometry, metric measure spaces, information geometry,
field theory, Lie algebra structure, sheaf theory, spectral theory,
stochastic analysis, categorical alignment, and unified synthesis.

%%% CHAPTER 2 %%%

\chapter[Latent Hierarchy and Constraint Reconstruction]{The Problem of Latent Hierarchy and Constraint Reconstruction}
\label{ch:hierarchy}

\section{Introduction}

Many scientific and cognitive domains confront the same fundamental
problem: a system possesses an underlying organization that is not
directly observable, while the available data consist only of partial
relations among the system's components.

\begin{definition}[Latent Hierarchy Reconstruction Problem]
Let $X$ be a finite set and let $T^\ast$ be an unknown rooted tree
whose leaves correspond to elements of $X$. Given relational
constraints $\mathcal{R}$ generated by $T^\ast$, the reconstruction
problem is to recover $\widehat{T}$ consistent with $\mathcal{R}$.
\end{definition}

\section{Trees and Hierarchical Structure}

\begin{definition}[Rooted Hierarchical Tree]
A rooted hierarchical tree on $X$ is a tree $T$ in which each leaf
corresponds uniquely to an element of $X$, each internal node
represents the union of its descendant leaves, and the root
corresponds to the entire set $X$.
\end{definition}

\begin{definition}[Laminar Family]
A family $\mathcal{F} \subseteq 2^X$ is \emph{laminar} if for every
pair $A,B \in \mathcal{F}$ one of the following holds:
$A \subseteq B$, $B \subseteq A$, or $A \cap B = \varnothing$.
\end{definition}

\section{Triplet Relations}

\begin{definition}[Rooted Triplet]
The rooted triplet $(ab|c)$ denotes the relation in which $a$ and $b$
share a more recent common ancestor with each other than either does
with $c$.
\end{definition}

\begin{definition}[Consistent Triplet System]
A triplet system $\mathcal{R}$ on $X$ is \emph{consistent} if there
exists a rooted tree $T$ such that $\mathcal{R} \subseteq
\mathcal{R}(T)$. Consistency is decidable in polynomial time via the
BUILD algorithm \cite{aho1981inferring}.
\end{definition}

\begin{figure}[htbp]
\centering
\begin{tikzpicture}[
  >=Latex, every node/.style={font=\small},
  leaf/.style={circle, fill, inner sep=1.4pt}
]
\coordinate (r) at (0,2.8);
\coordinate (u) at (-1.2,1.6);
\coordinate (v) at (1.2,1.6);
\coordinate (a) at (-2,0.4);
\coordinate (b) at (-0.6,0.4);
\coordinate (c) at (1.2,0.4);
\draw[thick] (r)--(u); \draw[thick] (r)--(v);
\draw[thick] (u)--(a); \draw[thick] (u)--(b);
\draw[thick] (v)--(c);
\node[leaf,label=below:$a$] at (a) {};
\node[leaf,label=below:$b$] at (b) {};
\node[leaf,label=below:$c$] at (c) {};
\node at (0,3.2) {$T^\ast$};
\node at (0,-0.4) {triplet witness: $(ab\mid c)$};
\draw[dashed] (-2.3,0.1) rectangle (-0.3,1.9);
\node at (-1.3,2.15) {$a,b$ closer};
\end{tikzpicture}
\caption{A rooted triplet $(ab\mid c)$ as the minimal witness of
branching structure. The pair $(a,b)$ shares a more recent common
ancestor, encoding a local constraint on the global tree topology.}
\label{fig:triplet-witness}
\end{figure}

\section{Reconstruction from Triplets}

\begin{theorem}[Triplet Reconstruction Theorem]
\label{thm:triplet-reconstruction}
Let $T$ be a rooted binary tree on $X$ with $|X| \geq 3$. If
$\mathcal{R}$ contains all triplets induced by $T$, then $T$ is the
unique tree whose topology is consistent with $\mathcal{R}$.
\end{theorem}

\begin{proof}
Each internal edge $e$ of $T$ induces a bipartition $X = L_e \sqcup
R_e$. Since $T$ is binary with $|X| \geq 3$, both sides contain at
least two leaves. For any $a,b \in L_e$ and $c \in R_e$, the triplet
$(ab|c)$ appears in $\mathcal{R}$. The bipartition is recovered from
$\mathcal{R}$ as the unique partition such that, for some $c$, the
triplet $(ab|c)$ holds for all $a,b$ on one side. Since the edge set
uniquely determines tree topology, $\mathcal{R}(T)$ uniquely
identifies $T$.
\end{proof}

\section{Constraint Operators and Entropy Reduction}

Each triplet observation $(ab|c)$ induces a constraint operator
\[
C_{(ab|c)} : 2^{\mathcal{H}} \rightarrow 2^{\mathcal{H}}, \qquad
C_{(ab|c)}(\mathcal{H}') =
\{ T \in \mathcal{H}' \mid T \text{ is consistent with } (ab|c) \}.
\]

\begin{proposition}
\label{prop:entropy-decrease}
Every triplet observation not already entailed by $\mathcal{H}_\mathcal{R}$
produces a strict reduction in the entropy of the hypothesis
distribution.
\end{proposition}

\begin{corollary}
If $\mathcal{R} = \mathcal{R}(T^\ast)$ is complete, the entropy
descent terminates at a delta distribution concentrated on $T^\ast$
in at most $|\mathcal{R}(T^\ast)|$ steps.
\end{corollary}

%%% CHAPTER 3 %%%

\chapter[Algorithmic Reconstruction]{Algorithmic Reconstruction and Computational Complexity}
\label{ch:algorithms}

\section{Introduction}

The Triplet Reconstruction Theorem provides a theoretical guarantee
of identifiability. This chapter examines how reconstruction is
carried out efficiently in practice.

\begin{figure}[htbp]
\centering
\begin{tikzpicture}[
  >=Latex, every node/.style={font=\small},
  box/.style={draw, rounded corners, minimum width=2.6cm,
    minimum height=0.95cm, align=center},
  arrow/.style={->, thick}
]
\node[box] (R) at (0,0) {$\mathcal{R}$\\triplet system};
\node[box] (G) at (3.8,0) {$G_{\mathcal{R}}$\\constraint graph};
\node[box] (CC) at (7.6,0) {connected\\components};
\node[box] (T1) at (11.4,1.3) {subtree $X_1$};
\node[box] (T2) at (11.4,0) {subtree $X_2$};
\node[box] (T3) at (11.4,-1.3) {subtree $X_3$};
\draw[arrow] (R)--(G); \draw[arrow] (G)--(CC);
\draw[arrow] (CC)--(T1); \draw[arrow] (CC)--(T2);
\draw[arrow] (CC)--(T3);
\node at (7,-2.2) {BUILD recursively reconstructs subtree structure};
\end{tikzpicture}
\caption{Schematic structure of the BUILD algorithm. A triplet system
$\mathcal{R}$ induces a constraint graph $G_{\mathcal{R}}$ whose
connected components determine recursive decomposition into subtrees.}
\label{fig:build-algorithm-schematic}
\end{figure}

\section{The BUILD Algorithm}

The BUILD algorithm \cite{aho1981inferring} constructs an undirected
graph $G_\mathcal{R}$ whose vertices are elements of $X$ with an edge
between $a$ and $b$ whenever some triplet $(ab|c) \in \mathcal{R}$.
Connected components correspond to clusters sharing a common ancestor.
The algorithm is applied recursively to each component with time
complexity $O(n^2)$ in the dense case.

\section{Incomplete and Probabilistic Systems}

When triplets are missing, one minimizes the violation count
$V(T) = |\{r \in \mathcal{R} \mid r \notin \mathcal{R}(T)\}|$,
which is NP-hard in general. The probabilistic approach treats
triplets as evidence: $p(T|\mathcal{R}) \propto p(T)\prod_r p(r|T)$,
corresponding to movement on the statistical manifold of tree
distributions toward lower entropy.

\section{Incremental Inference and Entropy Descent}

In incremental settings each arriving triplet reduces the hypothesis
space through a constraint operator application. When the triplet
system is complete, entropy reaches zero and the unique topology is
recovered, instantiating the abstract entropy descent of
\cref{ch:hierarchy} as a concrete algorithm.

%%% CHAPTER 4 %%%

\chapter[Ultrametric Geometry]{Ultrametric Geometry and the Metric Structure of Hierarchies}
\label{ch:ultrametric}

\section{Hierarchies as Metric Objects}

Every hierarchical tree induces a natural ultrametric on its leaves.
The space $\mathcal{U}(X)$ of all ultrametrics on a finite set $X$
with $|X| = n$ is a polyhedral cone in $\mathbb{R}^{\binom{n}{2}}$.

\begin{figure}[htbp]
\centering
\begin{tikzpicture}[every node/.style={font=\small}]
\draw[thick] (0,0) circle (2.2);
\draw[thick] (-0.8,0.1) circle (0.85);
\fill (-1.2,0.2) circle (1.6pt);
\fill (-0.5,-0.1) circle (1.6pt);
\fill (1.1,0.9) circle (1.6pt);
\node[below left] at (-1.2,0.2) {$x$};
\node[below] at (-0.5,-0.1) {$y$};
\node[right] at (1.1,0.9) {$z$};
\node at (-0.85,1.15) {$B_r(x)=B_r(y)$};
\node at (0,-2.7) {$d(x,y) < d(x,z)=d(y,z)$};
\end{tikzpicture}
\caption{Ultrametric ball geometry. The pair $(x,y)$ lies in a
smaller common ball excluding $z$, equivalent to the rooted triplet
relation $(xy\mid z)$.}
\label{fig:ultrametric-balls}
\end{figure}

\section{Ultrametric Spaces}

\begin{definition}[Ultrametric]
A metric $d$ is an \emph{ultrametric} if
\[
d(x,z) \le \max\{d(x,y), d(y,z)\}
\]
for all $x,y,z \in X$.
\end{definition}

\begin{proposition}
In any ultrametric space, at least two of the three pairwise
distances among any three points are equal and no smaller than the
third.
\end{proposition}

\section{Trees Induce Ultrametrics and Vice Versa}

For a rooted tree $T$ with nonneg\-ative edge lengths, define $d_T(x,y)$
as the height of the lowest common ancestor. Then $d_T$ is an
ultrametric. The converse is also true.

\begin{theorem}[Tree--Ultrametric Equivalence]
\label{thm:tree-ultrametric}
For any finite ultrametric $(X,d)$ there exists a rooted tree $T$
with $d = d_T$.
\end{theorem}

\begin{proof}
For each $r \geq 0$, the relation $x \sim_r y \Leftrightarrow d(x,y)
\leq r$ is an equivalence relation by the ultrametric inequality.
Equivalence classes nest as $r$ increases, producing a laminar family.
Each class becomes an internal node of $T$ at height $\inf\{r : x
\sim_r y\} = d(x,y)$, giving $d_T = d$.
\end{proof}

\section{Geometric Interpretation of Branching}

The triplet $(xy|z)$ corresponds to $d(x,y) < d(x,z) = d(y,z)$:
exactly the condition that $x$ and $y$ lie in a smaller ball
excluding $z$.

\section{Gromov--Hausdorff Convergence}

The Gromov--Hausdorff distance $d_{GH}(X,Y)$ is the infimum of
Hausdorff distances between isometric embeddings of $X$ and $Y$ into
a common ambient space. Ultrametric spaces form a closed subset of
the Gromov--Hausdorff metric space, so sequences of finite trees can
converge to continuous metric spaces in a rigorous sense.

%%% CHAPTER 5 %%%

\chapter[Limits of Hierarchies and Metric Measure Spaces]{Limits of Hierarchical Structures and Metric Measure Spaces}
\label{ch:metric-measure}

\section{Introduction}

Many applications require discrete hierarchical structures to be
understood as approximations to continuous geometric objects. This
chapter formalizes the transition using the theory of metric measure
spaces and Gromov--Hausdorff convergence, connecting the discrete
reconstruction theory to the continuous field models that follow.

\begin{figure}[htbp]
\centering
\begin{tikzpicture}[scale=1.1]
\draw (-4,0)--(-3,1)--(-2,0); \draw (-3,1)--(-3,2);
\node at (-3,-0.5) {$T_1$};
\draw (-1,0)--(0,1)--(1,0); \draw (0,1)--(0,2);
\draw (0,1)--(0.5,1.6); \draw (0,1)--(-0.5,1.6);
\node at (0,-0.5) {$T_2$};
\draw (2,0)--(3,1)--(4,0); \draw (3,1)--(3,2);
\draw (3,1)--(3.5,1.5); \draw (3,1)--(2.5,1.5);
\draw (3,1)--(3,1.7); \draw (3.5,1.5)--(3.7,2);
\node at (3,-0.5) {$T_n$};
\draw[->,thick] (4.6,0.7)--(5.8,0.7);
\draw[thick] (6.3,0) .. controls (7.1,1.5) and (8.1,1.5) .. (8.9,0);
\draw[thick] (7.6,1.5) .. controls (7.8,2.1) .. (7.6,2.7);
\node at (7.5,-0.5) {$T_\infty$};
\node at (2,3.1) {Gromov--Hausdorff convergence};
\end{tikzpicture}
\caption{A sequence of discrete hierarchical trees converging under
Gromov--Hausdorff limits to a continuum tree $T_\infty$.}
\label{fig:gh-convergence}
\end{figure}

\section{Metric Measure Spaces}

\begin{definition}[Metric Measure Space]
A \emph{metric measure space} is a triple $(X,d,\mu)$ consisting of
a compact metric space $(X,d)$ together with a Borel probability
measure $\mu$ on $X$.
\end{definition}

The measure $\mu$ represents the distribution of mass or probability
across the space, encoding relative weights of leaves or clusters.

\section{Gromov--Hausdorff Distance}

\begin{definition}[Hausdorff Distance]
For subsets $A,B$ of a metric space $(Z,d)$,
\[
d_H(A,B) = \max\Bigl(
\sup_{a \in A}\inf_{b \in B}d(a,b),\;
\sup_{b \in B}\inf_{a \in A}d(a,b)
\Bigr).
\]
\end{definition}

\begin{definition}[Gromov--Hausdorff Distance]
For compact metric spaces $(X,d_X)$ and $(Y,d_Y)$,
$d_{GH}(X,Y)$ is the infimum of Hausdorff distances between isometric
embeddings of $X$ and $Y$ into a common metric space.
\end{definition}

This makes the collection of compact metric spaces into a metric space
in its own right.

\section{Ultrametric Spaces are Closed under GH Convergence}

\begin{proposition}
The set of compact ultrametric spaces is closed under
Gromov--Hausdorff convergence.
\end{proposition}

\begin{proof}
Let $(X_n,d_n)$ converge to $(X,d)$ in $d_{GH}$. For $x,y,z \in X$
choose sequences $x_n,y_n,z_n$ converging to them. Since
$d_n(x_n,z_n) \leq \max\{d_n(x_n,y_n), d_n(y_n,z_n)\}$, taking
the limit gives $d(x,z) \leq \max\{d(x,y),d(y,z)\}$.
\end{proof}

The family of hierarchical metric spaces is therefore stable under
geometric limits.

\section{Scaling Limits of Trees}

\begin{definition}[Scaling Limit]
A metric space $(X,d)$ is a \emph{scaling limit} of $(X_n,d_n)$ if
$d_{GH}((X_n,d_n),(X,d)) \to 0$ after appropriate normalization.
\end{definition}

Scaling limits of large hierarchical trees often produce continuum
trees or fractal metric structures that appear in stochastic branching
process theory.

\section{Continuum Trees}

\begin{definition}[Continuum Tree]
A \emph{continuum tree} is a compact metric space $(T,d)$ in which
every pair of points is connected by a unique geodesic and the space
contains no cycles.
\end{definition}

Continuum trees arise as scaling limits of large random trees and
provide the continuous analogue of hierarchical structures.

\section{Entropy Measures on Metric Spaces}

The metric entropy at scale $\epsilon$ is
$H_\epsilon(X) = \log N_\epsilon(X)$,
where $N_\epsilon(X)$ is the minimum number of balls of radius
$\epsilon$ needed to cover $X$. The growth rate of $H_\epsilon(X)$
as $\epsilon \to 0$ reflects geometric complexity. For hierarchical
systems this reflects the diversity of structures present at
different scales, connecting the discrete entropy of
\cref{ch:information} to the geometry of metric measure spaces.

\section{From Trees to Field Configurations}

The Gromov--Hausdorff limit makes precise how discrete tree
reconstruction problems connect to continuous field inference. Triplet
consistency conditions become differential constraints on field
configurations in the limit, justifying the transition from
combinatorial constraint geometry to the RSVP field model.

%%% CHAPTER 6 %%%

\chapter{Geometry of Constraint Spaces}
\label{ch:constraint-geometry}

\section{Introduction}

The space of candidate hierarchical structures possesses a natural
polyhedral geometry. This chapter makes that geometry explicit and
reinterprets the reconstruction process as geometric descent within
the ultrametric cone.

\begin{figure}[htbp]
\centering
\begin{tikzpicture}[>=Latex, every node/.style={font=\small}]
\draw[thick] (0,0)--(4,0)--(2,3.4)--cycle;
\draw[thick] (2,0.9)--(3,1.7);
\draw[thick] (2,0.9)--(1,1.7);
\node at (2,-0.4) {$\mathcal{U}(X)$};
\node at (0.9,1.9) {$F_{T_1}$};
\node at (3.1,1.9) {$F_{T_2}$};
\node at (2,0.45) {$F_{T^\ast}$};
\fill (2.8,2.2) circle (1.6pt);
\fill (2.45,1.7) circle (1.6pt);
\fill (2.2,1.25) circle (1.6pt);
\draw[->,thick] (2.8,2.2)--(2.45,1.7);
\draw[->,thick] (2.45,1.7)--(2.2,1.25);
\draw[->,thick] (2.2,1.25)--(2.0,0.92);
\node[right] at (2.8,2.2) {$\mu_0$};
\node[below right] at (2.2,1.25) {$\mu_t$};
\end{tikzpicture}
\caption{Geometric reconstruction within the ultrametric cone.
Triplet constraints progressively restrict the admissible region,
driving the hypothesis distribution toward the face corresponding
to the true topology $T^\ast$.}
\label{fig:ultrametric-cone-descent}
\end{figure}

\section{The Ultrametric Cone}

\begin{proposition}
The set $\mathcal{U}(X)$ of all ultrametrics on $X$ is a convex
polyhedral cone in $\mathbb{R}^{\binom{n}{2}}$.
\end{proposition}

\section{Faces of the Cone and Tree Topologies}

\begin{proposition}
For each rooted tree topology $T$, the set of ultrametrics consistent
with $T$ forms a face of $\mathcal{U}(X)$.
\end{proposition}

Thus the combinatorial space of tree topologies appears geometrically
as the face lattice of the ultrametric cone.

\section{Constraint Intersections as Geometric Restrictions}

Each triplet $(ab|c)$ imposes the inequality
$d(a,b) < \min\{d(a,c), d(b,c)\}$, restricting the admissible region
to a subcone. The full constraint system $\mathcal{R}$ restricts to
the intersection of these subcones, which reduces to the face
corresponding to the true topology when $\mathcal{R}$ is complete.

\section{Entropy and Geometric Descent}

A probability measure $\mu$ on $\mathcal{U}(X)$ has entropy decreasing
as constraints restrict its support to smaller regions of the cone.
Inference is geometric descent in the polyhedral cone.

\section{Connection to Continuous Field Models}

This polyhedral geometry generalizes: in continuous field models the
constraint manifolds replace the cone faces, and differential
constraints replace the linear inequalities. The RSVP field framework
introduced in \cref{ch:rsvp} is the continuous analogue of this
structure.

%%% CHAPTER 7 %%%

\chapter[Entropy, Information Geometry, and Constraints]{Entropy, Information Geometry, and Constraint Propagation}
\label{ch:information}

\section{Shannon Entropy and Constraint Operators}

\begin{definition}[Shannon Entropy]
The entropy of a distribution $p$ over $\mathcal{H}$ is
$H(p) = -\sum_{h \in \mathcal{H}} p(h)\log p(h)$.
\end{definition}

\begin{definition}[Constraint Operator]
The constraint operator associated with observation $o$ is
\[
C_o : 2^{\mathcal{H}} \rightarrow 2^{\mathcal{H}}, \qquad
C_o(\mathcal{H}') = \{ h \in \mathcal{H}' \mid h \text{ consistent
with } o \}.
\]
\end{definition}

\begin{figure}[htbp]
\centering
\begin{tikzpicture}[>=Latex, every node/.style={font=\small}]
\draw[thick] (0,0) .. controls (2,1.8) and (5,1.8) .. (7,0);
\draw[thick] (0,0) .. controls (2,-1.8) and (5,-1.8) .. (7,0);
\draw[dashed] (1.1,0.55) .. controls (2.5,1.0) and (4.5,0.65) ..
  (6.0,0.1);
\draw[->,thick] (1.1,0.55)--(2.2,0.8);
\draw[->,thick] (2.8,0.84)--(3.8,0.72);
\draw[->,thick] (4.4,0.58)--(5.5,0.22);
\fill (1.1,0.55) circle (1.5pt);
\fill (6.0,0.1) circle (1.5pt);
\node[left] at (1.1,0.55) {$p_0$};
\node[right] at (6.0,0.1) {$\delta_{h^\ast}$};
\node at (3.5,-1.2) {statistical manifold of distributions};
\node at (3.6,1.35) {entropy descent};
\end{tikzpicture}
\caption{Inference as a trajectory on the statistical manifold.
Constraint accumulation drives the posterior from a diffuse
distribution $p_0$ toward concentration on the true structure.}
\label{fig:information-geometric-descent}
\end{figure}

\section{Bayesian Updating and Information Geometry}

Given prior $p(h)$ and observation $o$,
$p(h|o) = p(o|h)p(h)/\sum_{h'} p(o|h')p(h')$.
The Fisher information metric $g_{ij}(\theta) =
\mathbb{E}[\partial_i\log p \cdot \partial_j \log p]$ gives the
Riemannian structure on the statistical manifold; the KL divergence
$D_{KL}(p\|q) = \sum_h p(h)\log p(h)/q(h)$ approximates $g$
locally. Bayesian updating is geodesic flow under $g$ in the
exponential family.

\section{Variational Interpretation}

Inference minimizes a free energy functional
$F(p) = \mathbb{E}_p[E(h)] - \tau H(p)$,
where $E(h)$ counts constraint violations. In the tree reconstruction
setting $E(h)$ counts triplets in $\mathcal{R}$ that $h$ fails to
satisfy.

\section{Continuous Constraint Sets}

In continuous settings, observations impose constraint sets
$\mathcal{C}_i = \{(\Phi,\mathbf{v},S) \in \mathcal{X} \mid
\mathcal{O}_i(\Phi,\mathbf{v},S) \in [a_i,b_i]\}$, and inference
corresponds to projection onto $\bigcap_i \mathcal{C}_i$.

%%% CHAPTER 8 %%%

\chapter{The Relativistic Scalar--Vector Plenum}
\label{ch:rsvp}

\section{Field Variables and State Space}

\begin{definition}[RSVP Field State]
An RSVP field configuration is a triple $(\Phi, \mathbf{v}, S)$
where $\Phi : M \rightarrow \mathbb{R}$ is a scalar field,
$\mathbf{v} : M \rightarrow TM$ is a vector field, and
$S : M \rightarrow \mathbb{R}_{\ge 0}$ is an entropy density.
\end{definition}

\begin{definition}[RSVP State Space]
Fix $k \geq 2$. The RSVP state space is
$\mathcal{X} = H^k(M) \times H^k(M; TM) \times H^k(M)$.
\end{definition}

\begin{figure}[htbp]
\centering
\begin{tikzpicture}[>=Latex, every node/.style={font=\small},
  scale=1.0]
\draw[rounded corners, thick] (-4,-2.6) rectangle (4,2.6);
\draw[thin] (-3.2,-0.6) ellipse (1.3 and 0.7);
\draw[thin] (-3.2,-0.6) ellipse (0.8 and 0.4);
\draw[thin] (2.3,1.0) ellipse (1.2 and 0.6);
\draw[thin] (2.3,1.0) ellipse (0.7 and 0.35);
\node at (-3.2,-1.5) {$\Phi$ peak};
\node at (2.3,1.9) {$\Phi$ peak};
\draw[->,thick] (-2,-1)--(-1.2,-0.4);
\draw[->,thick] (-1,-1.3)--(0,-0.6);
\draw[->,thick] (0.2,-0.5)--(1.1,0.1);
\draw[->,thick] (1.3,-1.1)--(2.2,-0.5);
\draw[->,thick] (-0.6,1.0)--(0.2,1.6);
\draw[->,thick] (0.5,0.8)--(1.3,1.5);
\node at (3.5,-1.9) {$\mathbf{v}$};
\draw[dashed, thick] (-0.5,0.2) circle (1.4);
\node at (-0.5,1.8) {$S$ minimum};
\node at (0,-3.1) {RSVP field configuration on manifold $M$};
\end{tikzpicture}
\caption{Scalar salience field $\Phi$, conceptual flow $\mathbf{v}$,
and entropy density $S$ interacting across a representation manifold.
Peaks of $\Phi$ attract flow trajectories while entropy minima
stabilize emergent structures.}
\label{fig:rsvp-geometry}
\end{figure}

\section{Coupled Field Dynamics}

\begin{align}
\frac{\partial \Phi}{\partial t}
&= D_\Phi \nabla^2 \Phi
- \alpha_1 \nabla \cdot (\Phi \mathbf{v})
- \beta_1 S \Phi + \eta_\Phi, \label{eq:phi} \\
\frac{\partial \mathbf{v}}{\partial t}
&= D_v \nabla^2 \mathbf{v}
+ \alpha_2 \nabla \Phi
- \beta_2 S \mathbf{v} + \eta_v, \label{eq:v} \\
\frac{\partial S}{\partial t}
&= D_S \nabla^2 S
+ \gamma_1 |\nabla \cdot \mathbf{v}|
- \gamma_2 S + \eta_S. \label{eq:S}
\end{align}

\begin{remark}
The advection term $-\alpha_1 \nabla \cdot (\Phi \mathbf{v})$ models
conservative transport and is well-typed. The term $-\gamma_2 S$
ensures entropy dissipates toward zero in the absence of flow
divergence. The full dynamics are a dissipative perturbation of
gradient flow on the energy functional $\mathcal{E}$ below.
\end{remark}

\section{Energy Functional}

\begin{definition}[RSVP Energy]
$\mathcal{E}(\Phi,\mathbf{v},S) =
\int_M \bigl(\tfrac{1}{2}|\nabla\Phi|^2
+ \tfrac{1}{2}|\mathbf{v}|^2 + \lambda S\bigr)\,dV$.
\end{definition}

\begin{proposition}
\label{prop:energy-convex}
$\mathcal{E}$ is strictly convex in $\Phi$ and $\mathbf{v}$
separately and affine in $S$. When coupling terms vanish it has a
unique minimizer for each boundary condition. Emergent structures
arise in the coupled regime where convexity is broken.
\end{proposition}

%%% CHAPTER 9 %%%

\chapter[Well-Posedness of the RSVP Equations]{Well-Posedness and Dynamical Regimes of the RSVP Field Equations}
\label{ch:wellposedness}

\section{Local Existence}

\begin{theorem}[Local Existence]
\label{thm:local-existence}
Let $(\Phi_0,\mathbf{v}_0,S_0) \in \mathcal{X}$ with $k \geq 2$ and
suppose the stochastic forcing terms are bounded in $H^k$. Then there
exists $T > 0$ and a unique solution $(\Phi(t),\mathbf{v}(t),S(t))$
of \eqref{eq:phi}--\eqref{eq:S} on $[0,T]$.
\end{theorem}

\begin{proof}
The diffusion terms generate analytic semigroups on $H^k(M)$. The
nonlinear coupling terms are locally Lipschitz in $\mathcal{X}$ by the
Sobolev multiplication theorem for $k > \dim M/2$. The Picard
iteration method applied to the integral formulation yields a unique
local solution by the Banach fixed-point theorem.
\end{proof}

\section{Energy Dissipation and Dynamical Regimes}

\begin{proposition}
If stochastic forcing vanishes, $d\mathcal{E}/dt \leq 0$ for
sufficiently small coupling parameters.
\end{proposition}

\begin{figure}[htbp]
\centering
\begin{tikzpicture}[>=Latex, every node/.style={font=\small}]
\draw[->,thick] (0,0)--(6.2,0) node[right] {coupling strength};
\draw[->,thick] (0,0)--(0,4.4) node[above] {diffusion / damping};
\draw[thick] (0.8,3.5)--(5.3,0.8);
\draw[thick] (1.2,3.9)--(5.7,1.2);
\node at (1.6,3.0) {diffusive regime};
\node at (4.8,3.0) {critical / transient};
\node at (4.6,1.0) {structured regime};
\fill (4.6,1.0) circle (1.5pt);
\draw[->,thick] (3.2,2.0)--(4.4,1.15);
\end{tikzpicture}
\caption{Qualitative dynamical regimes of the RSVP equations. Strong
diffusion suppresses structure; strong coupling supports pattern
formation and attractor dynamics.}
\label{fig:rsvp-regimes}
\end{figure}

Three qualitatively distinct regimes arise. When diffusion dominates,
fields relax to uniform configurations. When coupling is sufficiently
strong, localized patterns emerge through self-organization. An
intermediate critical regime produces transient patterns. The entropy
field mediates transitions between these regimes through
\eqref{eq:S}.

%%% CHAPTER 10 %%%

\chapter[Amplitwist Operators]{Amplitwist Operators and Local Geometric Transformations}
\label{ch:amplitwist}

\section{Amplitwists in Complex Analysis}

For a holomorphic function $f : \mathbb{C} \rightarrow \mathbb{C}$,
the derivative $f'(z_0) = \rho e^{i\theta}$ acts locally as
$v \mapsto \rho R(\theta) v$: rotation by $\theta$ followed by
scaling by $\rho$ \cite{needham1997}.

\begin{definition}[Amplitwist Operator]
An amplitwist operator on $\mathbb{R}^2$ is
$A = s R(\theta)$ with $s \in \mathbb{R}_{>0}$ and $R(\theta)$ the
rotation matrix. For $n=2$, the group of amplitwists is
$\mathbb{R}_{>0} \times SO(2) \cong \mathbb{R}_{>0} \times U(1)$,
the multiplicative group of nonzero complex numbers.
\end{definition}

\begin{figure}[htbp]
\centering
\begin{tikzpicture}[scale=1.1, >=Latex, every node/.style={font=\small}]
\draw[->] (-3,0)--(3,0);
\draw[->] (0,-3)--(0,3);
\draw[thick] (0,0) circle (1);
\draw[thick,rotate=40] (0,0) ellipse (1.6 and 0.7);
\draw[->,thick] (0,0)--(1,0);
\draw[->,thick] (0,0)--({1.6*cos(40)},{1.6*sin(40)});
\node at (1.2,-0.3) {$v$};
\node at (1.2,1.1) {$A_x v$};
\node at (0,-3.4) {$A_x = \alpha_x R_x$};
\end{tikzpicture}
\caption{Local geometric action of an amplitwist operator.
Infinitesimal neighborhoods undergo rotation and scaling in the
tangent space, generalizing the complex-analytic derivative to
representation manifolds.}
\label{fig:amplitwist-differential}
\end{figure}

\section{Generalized Amplitwists and Realization Conditions}

\begin{definition}[Generalized Amplitwist]
A generalized amplitwist at $x \in M$ is
$A_x = \alpha_x R_x : T_x M \rightarrow T_x M$ with $\alpha_x > 0$
and $R_x \in SO(n)$.
\end{definition}

\begin{definition}[Amplitwist Realization]
\label{def:realization}
A system $\mathcal{S}$ with map $F : \Sigma \rightarrow \Sigma$
\emph{realizes} $A_x$ with tolerance $\varepsilon$ if there exists
a measurable encoding map $\psi : \Sigma \rightarrow T_x M$ with
$\sup_{u \in \mathcal{U}_x} \|\psi(F(u)) - A_x\psi(u)\| \leq
\varepsilon$.
\end{definition}

\section{Local Field Transformations and Entropy Regulation}

Application of $A_x = \alpha_x R_x$ at $x$ produces perturbations
$\delta\Phi_x = (\alpha_x-1)\Phi(x)\chi_x$ and
$\delta\mathbf{v}_x = (R_x-I)\mathbf{v}(x)\chi_x$.

\begin{definition}[Entropy-Modulated Amplitwist]
$A^{(\beta)}_x = \alpha_x e^{-\beta S(x)} R_x$.
\end{definition}

\begin{proposition}
\label{prop:entropy-modulation}
$\|A^{(\beta)}_x\| = \alpha_x e^{-\beta S(x)} \leq \alpha_x$ with
equality iff $S(x)=0$, and $\lim_{\beta\to\infty}\|A^{(\beta)}_x\|
= 0$ for any $S(x)>0$.
\end{proposition}

\section{Cortical Columns as Transformation Units}

Cortical columns \cite{mountcastle1957} realize amplitwist operators
in the sense of \cref{def:realization} when population recordings
admit an encoding map $\psi$ with bounded approximation error.
Predicted signatures: rotation-like trajectories in neural population
codes \cite{dicarlo2015}, and entropy--amplitude anti-correlation
\cite{richman2000,higgins2017}.

%%% CHAPTER 11 %%%

\chapter[Lie Algebra Structure of Cascades]{Lie Algebra Structure of Amplitwist Cascades}
\label{ch:amplitwist-lie}

\section{Introduction}

The amplitwist transformations of \cref{ch:amplitwist} form a Lie
group, and recursive cascade composition is governed by the algebraic
structure of its Lie algebra. This chapter develops that structure
explicitly and connects cascade dynamics to geometric control theory.

\section{The Amplitwist Lie Group}

\begin{proposition}
The set of amplitwist transformations forms a Lie group isomorphic to
$G = \mathbb{R}_{>0} \times SO(n)$.
\end{proposition}

\begin{proof}
The amplitude scalings $\mathbb{R}_{>0}$ and rotations $SO(n)$ each
form Lie groups under their respective operations. Their direct
product $G = \mathbb{R}_{>0} \times SO(n)$ inherits a Lie group
structure with multiplication
$(\alpha_1,R_1)(\alpha_2,R_2) = (\alpha_1\alpha_2, R_1 R_2)$.
Each amplitwist corresponds uniquely to such a pair.
\end{proof}

\section{The Amplitwist Lie Algebra}

\begin{definition}[Amplitwist Generator]
An \emph{amplitwist generator} is an element $X = (\lambda, T)$
where $\lambda \in \mathbb{R}$ represents infinitesimal amplitude
change and $T \in \mathfrak{so}(n)$ is a skew-symmetric matrix
representing an infinitesimal rotation.
\end{definition}

The Lie algebra is $\mathfrak{g} = \mathbb{R} \oplus \mathfrak{so}(n)$.
The Lie bracket satisfies
\[
[(\lambda_1, T_1), (\lambda_2, T_2)] = (0, [T_1, T_2]),
\]
so amplitude scalings commute with everything, while rotational
generators interact through the standard bracket of $\mathfrak{so}(n)$.

\section{Cascade Composition and BCH Expansion}

\begin{theorem}[Baker--Campbell--Hausdorff]
For generators $X,Y \in \mathfrak{g}$,
\[
\exp(X)\exp(Y) = \exp\!\Bigl(X + Y + \tfrac{1}{2}[X,Y]
+ \tfrac{1}{12}[X,[X,Y]] - \tfrac{1}{12}[Y,[X,Y]] + \cdots\Bigr).
\]
\end{theorem}

\begin{figure}[htbp]
\centering
\begin{tikzpicture}[node distance=2.2cm, >=Stealth,
  every node/.style={font=\small}]
\node (x0) {$x_0$};
\node[right of=x0] (x1) {$x_1$};
\node[right of=x1] (x2) {$x_2$};
\node[right of=x2] (x3) {$x_3$};
\draw[->] (x0) -- node[above] {$e^{X_1}$} (x1);
\draw[->] (x1) -- node[above] {$e^{X_2}$} (x2);
\draw[->] (x2) -- node[above] {$e^{X_3}$} (x3);
\node at (3.3,-0.9) {cascade flow on $\mathbb{R}_{>0}\times SO(n)$};
\end{tikzpicture}
\caption{Layered amplitwist cascade as successive exponentials of Lie
algebra generators. BCH corrections accumulate torsion through Lie
brackets $[X_i, X_j]$.}
\label{fig:lie-cascade}
\end{figure}

For a cascade of $N$ layers, the effective generator is
\[
X_{\mathrm{eff}} = \sum_{k=1}^N X_k
+ \sum_{i<j} \tfrac{1}{2}[X_i,X_j] + O(\epsilon^3),
\]
where the torsion term $\sum_{i<j}\tfrac{1}{2}[T_i,T_j]$ captures
the non-commutativity of successive rotations and explains the
curvature observed in linguistic and cultural cascade dynamics.

\section{Cascade Flow and Continuous Limit}

\begin{definition}[Cascade Flow]
A \emph{cascade flow} is a path $A(t) = \exp(X(t))$ in $G$
generated by a time-dependent generator $X(t) \in \mathfrak{g}$.
\end{definition}

The representation state evolves as $\dot{x}(t) = X(t) \cdot x(t)$,
connecting amplitwist dynamics to the classical theory of flows on
Lie groups.

\section{Energy Constraints and Geometric Control}

A transformation $A$ is energetically admissible when
$\langle \nabla_\mathcal{X}\mathcal{E}, \delta_A \rangle \leq 0$.
Infinitesimally, admissible generators lie in the subspace of
$\mathfrak{g}$ for which $\frac{d}{dt}\mathcal{E}(x(t)) \leq 0$.
This formulation identifies the cascade system as a control system
$\dot{x} = \sum_i u_i(t)X_i(x)$ on the representation manifold,
where the generators $X_i$ represent elementary transformations and
the cascade weights $u_i(t)$ are control inputs.

\section{Cascade Stability via Lie Algebra Norms}

\begin{theorem}[Cascade Stability]
\label{thm:cascade-stability}
Let $C = \max_j\|T_j\|_{\mathrm{op}}$ and $\epsilon_\mathrm{crit}
= 1/(CN)$. If $\epsilon_j < \epsilon_\mathrm{crit}$ for all
$j \leq N$, then $\|\mathcal{R}_k - \mathrm{id}\|_{C^1} < 1$ for
all $k \leq N$ and the vorticity satisfies
\[
\lim_{N\to\infty} \xi^{(N)} \leq
\frac{C}{\mathrm{Vol}(M)}\int_M \|\nabla \times T_N(x)\|\,dx < \infty.
\]
\end{theorem}

\begin{remark}
The threshold $\epsilon_\mathrm{crit} = 1/(CN)$ is sufficient but
not tight. A sharp characterization in terms of spectral properties
of $\{T_j\}$ remains open.
\end{remark}

%%% CHAPTER 12 %%%

\chapter[Recursive Cascades and Semantic Dynamics]{Recursive Amplitwist Cascades and Multiscale Semantic Dynamics}
\label{ch:cascades}

\section{Layered Deformation}

Let $M$ be a smooth compact manifold. The $k$-layer deformation is
\[
\mathcal{R}_k(x) = x + \sum_{j=1}^k \epsilon_j T_j(x), \quad
T_j \in \mathfrak{so}(n).
\]
This is the first-order approximation to $\exp(\epsilon_1 T_1)\circ
\cdots \circ \exp(\epsilon_k T_k)$; BCH corrections introduce
torsion $\Theta^{(N)} = \sum_{i<j}\epsilon_i\epsilon_j[T_i,T_j]
+ O(\epsilon^3)$ as shown in \cref{ch:amplitwist-lie}.

The layer-$k$ amplitwist is
\[
A^{(k)}(x) = w_k(x)\cdot A(\mathcal{R}_k(x)), \qquad
w_k(x) = \exp(-\lambda S(x)).
\]

\begin{figure}[htbp]
\centering
\begin{tikzpicture}[>=Latex, every node/.style={font=\small},
  layer/.style={draw, circle, minimum size=1.3cm, align=center}]
\node[layer] (R1) at (0,0) {$\mathcal{R}_1$};
\node[layer] (R2) at (3,0) {$\mathcal{R}_2$};
\node[layer] (R3) at (6,0) {$\mathcal{R}_3$};
\node[layer] (A) at (9,0) {attractor};
\draw[->,thick] (R1)--(R2);
\draw[->,thick] (R2)--(R3);
\draw[->,thick] (R3)--(A);
\draw[decorate,decoration={snake,amplitude=1.2pt,segment length=8pt},thick]
  (0.6,0.55)--(2.4,0.55);
\draw[decorate,decoration={snake,amplitude=1.2pt,segment length=8pt},thick]
  (3.6,0.55)--(5.4,0.55);
\draw[decorate,decoration={snake,amplitude=1.2pt,segment length=8pt},thick]
  (6.6,0.55)--(8.4,0.55);
\node at (1.5,1.0) {$\epsilon_1 T_1$};
\node at (4.5,1.0) {$\epsilon_2 T_2$};
\node at (7.5,1.0) {$\epsilon_3 T_3$};
\node at (4.5,-1.3) {recursive deformation across scales};
\end{tikzpicture}
\caption{Recursive amplitwist cascade. Deformation layers
$\mathcal{R}_k$ apply Lie-algebraic perturbations whose cumulative
effect generates semantic attractors.}
\label{fig:amplitwist-cascade}
\end{figure}

\section{Multiscale Structure and Renormalization Analogy}

Define the effective spatial scale of layer $k$ as
$\ell_k \sim \|\epsilon_k T_k\|^{-1}$. Early layers with large
$\epsilon_k$ produce large-scale deformations; later layers refine
local details. This parallels Wilsonian renormalization.

\begin{figure}[htbp]
\centering
\begin{tikzpicture}[scale=1.1]
\draw[thick] (0,0) circle (2); \node at (0,2.5) {coarse scale};
\draw[thick] (5,0) circle (1.4); \node at (5,2.5) {intermediate};
\draw[thick] (9,0) circle (0.8); \node at (9,2.5) {fine scale};
\draw[->,thick] (2.2,0)--(3.7,0);
\draw[->,thick] (6.5,0)--(8.1,0);
\node at (3.3,-0.8) {$\ell_1$};
\node at (7.1,-0.8) {$\ell_2$};
\end{tikzpicture}
\caption{Cascade deformation across scales. Early layers generate
large-scale transformations while later layers refine local structure.}
\label{fig:renorm-scales}
\end{figure}

\section{Applications to Linguistic Evolution}

The cascade models linguistic change as:
$\mathcal{R}_1$: phonetic drift ($T_1$ = vowel shift);
$\mathcal{R}_2$: grammaticalization ($T_2$ = aspect-to-tense);
$\mathcal{R}_3$: semantic bleaching ($T_3$ = metaphor decay).
The torsion $[T_1,T_2]$ captures documented interactions between
sound change and grammaticalization. Empirical testing is tractable
via diachronic word embedding corpora.

%%% CHAPTER 13 %%%

\chapter{Sheaf-Theoretic Constraint Propagation}
\label{ch:sheaf-constraints}

\section{Introduction}

Triplet reconstruction, ultrametric geometry, and RSVP field
inference all share a common structure: local information must
propagate coherently to produce globally consistent configurations.
Sheaf theory formalizes this phenomenon precisely.

\section{Presheaves and Sheaves}

\begin{definition}[Presheaf]
A \emph{presheaf} $\mathcal{F}$ on a topological space $M$ assigns
to each open set $U \subseteq M$ a set $\mathcal{F}(U)$ (the
\emph{sections over $U$}) together with restriction maps
$\rho_{UV} : \mathcal{F}(U) \rightarrow \mathcal{F}(V)$ for
$V \subseteq U$, satisfying $\rho_{UU} = \mathrm{id}$ and
$\rho_{VW} \circ \rho_{UV} = \rho_{UW}$.
\end{definition}

\begin{definition}[Sheaf]
A presheaf $\mathcal{F}$ is a \emph{sheaf} if for every open cover
$\{U_i\}$ of $U$: (i) sections agreeing on all $U_i$ are equal;
(ii) locally compatible sections $s_i \in \mathcal{F}(U_i)$ with
$\rho_{U_i,U_{ij}}(s_i) = \rho_{U_j,U_{ij}}(s_j)$ on all overlaps
$U_{ij} = U_i \cap U_j$ glue to a unique section $s \in
\mathcal{F}(U)$.
\end{definition}

\section{Constraint Systems as Sheaves}

For each open set $U \subseteq M$, define
\[
\mathcal{F}(U) = \{ x \in \mathcal{X} \mid x \text{ satisfies
all constraints in } U \}.
\]

\begin{proposition}
The assignment $U \mapsto \mathcal{F}(U)$ defines a presheaf on $M$.
When the constraint system is locally consistent, it extends to a
sheaf.
\end{proposition}

\begin{definition}[Global Section]
A \emph{global section} of $\mathcal{F}$ is an element of
$\mathcal{F}(M)$, i.e., a configuration satisfying all constraints
simultaneously.
\end{definition}

\begin{proposition}
If the constraint system is consistent, the constraint presheaf
admits at least one global section.
\end{proposition}

\section{Sheaf Interpretation of Triplet Reconstruction}

Each triplet constraint restricts tree topologies on a small subset
of leaves. Local compatibility ensures partial tree structures agree
on overlapping leaf sets. When the triplet system is consistent,
these local structures glue into a single global tree topology: the
Triplet Reconstruction Theorem (\cref{thm:triplet-reconstruction})
is the existence of a global section of the constraint sheaf over the
leaf set.

\section{Sheaf Interpretation of RSVP Inference}

For each open region $U \subseteq M$, define
$\mathcal{F}(U) = \{(\Phi,\mathbf{v},S) \mid (\Phi,\mathbf{v},S)$
satisfies field constraints in $U\}$.
Compatible field solutions on overlapping regions assemble into a
global field configuration. RSVP inference is therefore a sheaf
gluing problem, and well-posedness of the field equations
(\cref{thm:local-existence}) guarantees the existence of local
sections.

\section{Cohomological Obstructions}

When constraints conflict, global sections may fail to exist despite
apparent local consistency. In sheaf theory, such failures are
measured by elements of the first cohomology group $H^1(M,\mathcal{F})$.
Non-trivial cohomology classes correspond to irresolvable
inconsistencies in the constraint network. This connects the
framework to applied topology and sheaf-based data analysis.

\section{Relation to Categorical Alignment}

Sheaves form categories whose morphisms are maps preserving
restriction structure. Constraint propagation corresponds to
functorial transport of local information. This connects directly to
the alignment framework of \cref{ch:alignment}: the action functor
$G : \mathcal{M} \rightarrow \mathcal{A}$ is required to be a
morphism of sheaves, not merely a map of underlying sets.

%%% CHAPTER 14 %%%

\chapter{Spectral Geometry of Constraint Operators}
\label{ch:spectral-constraints}

\section{Introduction}

Constraint systems can be interpreted as linear operators on spaces
of configurations. Spectral analysis of these operators reveals the
rigidity of the constraint system, the rate of inference convergence,
and the degree to which local constraints determine global structure.

\section{Constraint Laplacians}

Let $\{o_i\}_{i=1}^m$ be a set of constraints with associated
operators $C_{o_i}$. The \emph{constraint Laplacian} is
\[
L = \sum_{i=1}^m (I - C_{o_i}).
\]

\begin{proposition}
Configurations satisfying all constraints lie in the kernel of $L$.
\end{proposition}

\begin{proof}
If $Lv = 0$, then $(I - C_{o_i})v = 0$ for every $i$, so
$C_{o_i}v = v$: the configuration satisfies every constraint.
\end{proof}

\section{Spectral Gap and Reconstruction Stability}

Let $\lambda_2$ denote the smallest nonzero eigenvalue of $L$.

\begin{definition}[Spectral Gap]
The \emph{spectral gap} is $\gamma = \lambda_2$.
\end{definition}

\begin{proposition}
A large spectral gap implies unique, rapidly convergent inference.
A small spectral gap indicates that multiple configurations nearly
satisfy the constraints, producing reconstruction ambiguity.
\end{proposition}

This connects the reconstruction theory to the Cheeger isoperimetric
inequality and spectral graph theory via the Cheeger bound
$\gamma \geq h^2/2$, where $h$ is the edge expansion constant of the
constraint graph.

\section{Diffusion Interpretation}

The equation $\frac{dp}{dt} = -Lp$ defines a diffusion process on
the hypothesis space.

\begin{proposition}
The diffusion dynamics converge to the kernel of $L$ exponentially
at rate $\gamma$.
\end{proposition}

Thus constraint satisfaction is the equilibrium state of a diffusion
process, and the spectral gap governs the rate of entropy descent.

\section{Continuous Limit and the Laplace--Beltrami Operator}

When the hypothesis space becomes continuous, the constraint
Laplacian generalizes to the Laplace--Beltrami operator on the
configuration manifold $M$. Eigenfunctions describe intrinsic modes
of variation. Low-frequency modes correspond to global structural
features; high-frequency modes correspond to local fluctuations. This
spectral decomposition provides a geometric description of how
constraints shape the representation manifold and complements the
RSVP field equations \eqref{eq:phi}--\eqref{eq:S}.

\section{Connection to Entropy Descent}

Under the diffusion $\frac{dp}{dt} = -Lp$, entropy $H(p)$ decreases
monotonically at rate bounded below by $\gamma \cdot H(p)$. Thus the
spectral properties of the constraint operator directly govern the
rate of entropy reduction, closing the connection between the spectral
framework and the variational analysis of \cref{ch:information}.

%%% CHAPTER 15 %%%

\chapter[Stochastic Extensions of RSVP Dynamics]{Stochastic Extensions of RSVP Field Dynamics}
\label{ch:stochastic-rsvp}

\section{Introduction}

The deterministic RSVP dynamics of \cref{ch:rsvp} model the
evolution of $(\Phi,\mathbf{v},S)$ under coupled partial differential
equations. Many real systems evolve under stochastic perturbations.
This chapter extends the RSVP framework to a stochastic setting via
the theory of stochastic partial differential equations.

\section{Stochastic Field Equations}

The stochastic RSVP equations are
\begin{align}
\frac{\partial\Phi}{\partial t}
&= D_\Phi\nabla^2\Phi - \alpha_1\nabla\cdot(\Phi\mathbf{v})
 - \beta_1 S\Phi + \sigma_\Phi\,\dot{W}_\Phi, \\
\frac{\partial\mathbf{v}}{\partial t}
&= D_v\nabla^2\mathbf{v} + \alpha_2\nabla\Phi - \beta_2 S\mathbf{v}
 + \sigma_v\,\dot{W}_v, \\
\frac{\partial S}{\partial t}
&= D_S\nabla^2 S + \gamma_1|\nabla\cdot\mathbf{v}| - \gamma_2 S
 + \sigma_S\,\dot{W}_S,
\end{align}
where $\dot{W}_\Phi$, $\dot{W}_v$, $\dot{W}_S$ are spatially
correlated Gaussian white noise fields.

\begin{definition}[Gaussian Noise]
A stochastic field $\eta(x,t)$ is Gaussian white noise if
$\mathbb{E}[\eta(x,t)] = 0$ and
$\mathbb{E}[\eta(x,t)\eta(x',t')] = \sigma^2\delta(x-x')\delta(t-t')$.
\end{definition}

\section{Fokker--Planck Equation}

The stochastic RSVP system defines a Markov process on $\mathcal{X}$.
The probability density $P(\Phi,\mathbf{v},S,t)$ over field
configurations evolves via the Fokker--Planck equation
\[
\frac{\partial P}{\partial t}
= -\nabla\cdot(FP) + D\nabla^2 P,
\]
where $F$ is the deterministic drift from \eqref{eq:phi}--\eqref{eq:S}
and $D = \tfrac{1}{2}(\sigma_\Phi^2 + \sigma_v^2 + \sigma_S^2)$.

\section{Stationary Distributions and Free Energy}

\begin{definition}[Stationary Distribution]
$P^\ast$ is stationary if $\frac{\partial P^\ast}{\partial t} = 0$.
\end{definition}

Under detailed balance, the stationary distribution takes the Gibbs
form $P^\ast(x) \propto \exp(-\mathcal{E}(x)/\tau)$ where $\tau$ is
the effective temperature and $\mathcal{E}$ is the RSVP energy
functional. The stochastic dynamics therefore minimize the free energy
functional $F[P] = \mathbb{E}_P[\mathcal{E}] - \tau H(P)$, directly
connecting the statistical physics formulation to the variational
inference framework of \cref{ch:information}.

\section{Fluctuation-Dissipation Relation}

In physical equilibrium systems the noise amplitudes $\sigma_i$ and
the damping coefficients $\beta_i$ satisfy the
fluctuation-dissipation relation $\sigma_i^2 = 2\beta_i\tau$,
ensuring convergence to the correct Gibbs stationary distribution.
Departing from this relation introduces non-equilibrium dynamics,
corresponding to driven systems that maintain persistent flows.

\section{Stochastic Cascades}

Amplitwist generators become stochastic: $X_k(t) = X_k +
\xi_k(t)$ where $\xi_k$ is a random perturbation. The resulting
cascade is a stochastic flow on the Lie group $G$, governed by a
stochastic differential equation on $\mathbb{R}_{>0} \times SO(n)$.
This provides a model of noisy hierarchical inference processes in
which exploration and exploitation are balanced through the noise
amplitude.

\section{Implications for Inference and Alignment}

Stochastic field dynamics allow the system to escape local minima of
$\mathcal{E}$ and explore the configuration space more effectively.
For alignment, stochasticity introduces additional complexity: the
action functor $G$ must preserve invariants not only for
deterministic trajectories but in expectation over stochastic
realizations, requiring an averaged alignment condition
$\mathbb{E}[\alpha(G,I)] = 1$ that is strictly stronger than
pointwise alignment.

%%% CHAPTER 16 %%%

\chapter[Worked Examples of Cascade Dynamics]{Worked Examples of Amplitwist Cascade Dynamics}
\label{ch:examples}

\section{Introduction}

This chapter illustrates the cascade mechanisms of
\cref{ch:amplitwist-lie,ch:cascades} through concrete
two-dimensional examples, making visible how local transformations
accumulate to produce stable global structures.

\begin{figure}[htbp]
\centering
\begin{tikzpicture}[scale=0.95, >=Latex, every node/.style={font=\small}]
\draw[thin,->] (-2.5,0)--(2.5,0);
\draw[thin,->] (0,-2.5)--(0,2.5);
\draw[->,thick] (-1.8,0)--(-1.1,0);
\draw[->,thick] (1.1,0)--(1.8,0);
\draw[->,thick] (0,1.8)--(0,1.1);
\draw[->,thick] (0,-1.1)--(0,-1.8);
\draw[->,thick,rotate around={30:(0,0)}] (-1.8,0)--(-1.1,0);
\draw[->,thick,rotate around={30:(0,0)}] (1.1,0)--(1.8,0);
\draw[->,thick,rotate around={30:(0,0)}] (0,1.8)--(0,1.1);
\draw[->,thick,rotate around={30:(0,0)}] (0,-1.1)--(0,-1.8);
\node at (0,2.9) {original and transformed local flow};
\node at (0,-2.9) {$\mathbf{v}_1 = A\mathbf{v}_0$};
\end{tikzpicture}
\caption{Worked example of amplitwist action on a local vector field.
The original flow pattern is rotated and rescaled under the local
operator $A$.}
\label{fig:worked-flow-example}
\end{figure}

\section{Two-Layer Cascade}

For the initial saddle field $\mathbf{v}_0(x,y) = (-x,y)^T$, two
amplitwists $A_1 = s_1 R(\theta_1)$ and $A_2 = s_2 R(\theta_2)$
compose as
$A_2 A_1 = (s_1 s_2) R(\theta_1+\theta_2)$,
demonstrating that complex global patterns emerge from simple
compositional operations.

\section{Entropy-Modulated Cascade and Attractor Formation}

With entropy field $S$ concentrated near $x^\ast$, the
entropy-modulated operator $A^{(\beta)}(x,y) = \alpha e^{-\beta
S(x,y)} R(\theta)$ focuses transformations on the entropy minimum.

\begin{figure}[htbp]
\centering
\begin{tikzpicture}[scale=1.0, >=Latex, every node/.style={font=\small}]
\draw[thick,domain=0:4.5*pi,samples=120]
  plot ({0.06*\x*cos(\x r)},{0.06*\x*sin(\x r)});
\draw[thick,domain=0:4.5*pi,samples=120]
  plot ({0.06*\x*cos(\x r+1.1)},{0.06*\x*sin(\x r+1.1)});
\fill (0,0) circle (2.2pt);
\node at (0,-2.2) {cascade attractor};
\end{tikzpicture}
\caption{Conceptual attractor produced by recursive entropy-modulated
cascades. Trajectories spiral toward a stable semantic configuration.}
\label{fig:cascade-attractor}
\end{figure}

\section{Cascade Manifold Deformation}

\begin{figure}[htbp]
\centering
\begin{tikzpicture}[scale=1.0, >=Latex, every node/.style={font=\small}]
\draw[thick] (-4,-1.5) .. controls (-2,1.5) and (2,1.5) .. (4,-1.5);
\node at (0,-2) {$M$};
\draw[thick,dashed] (-4,-0.8)..controls(-2,2.0)and(2,2.0)..(4,-0.8);
\draw[thick,dotted] (-4,-0.1)..controls(-2,2.4)and(2,2.4)..(4,-0.1);
\draw[->,thick] (4.3,-1.4)--(4.3,-0.7);
\draw[->,thick] (4.3,-0.7)--(4.3,0.0);
\node[right] at (4.3,-1.4) {$\mathcal{R}_1$};
\node[right] at (4.3,-0.7) {$\mathcal{R}_2$};
\node[right] at (4.3,0.0) {$\mathcal{R}_3$};
\node at (0,2.7) {recursive manifold deformation};
\end{tikzpicture}
\caption{Recursive deformation of a representation manifold under
amplitwist cascades. Each layer $\mathcal{R}_k$ applies a controlled
Lie-algebraic perturbation.}
\label{fig:cascade-deformation}
\end{figure}

\section{Discrete Network Example}

When $M$ consists of a finite network of conceptual states, an
amplitwist modifies transition weights (amplification) and transition
directions (rotation). Successive cascade layers reinforce certain
pathways until the network converges to a stable configuration,
corresponding to the attractor predicted by
\cref{thm:cascade-stability}.

%%% CHAPTER 17 %%%

\chapter[Empirical Observability]{Empirical Observability and Measurement of Field Structures}
\label{ch:observability}

\section{Introduction}

For the framework to have empirical significance, the abstract
quantities of the theory must correspond to measurable features.
Let $\Sigma$ denote observable signals, $M$ the latent
representation manifold, and $\psi : \Sigma \rightarrow M$ an
encoding map.

\begin{figure}[htbp]
\centering
\begin{tikzpicture}[>=Latex, every node/.style={font=\small},
  box/.style={draw, rounded corners, minimum width=2.9cm,
    minimum height=1cm, align=center},
  arrow/.style={->, thick}]
\node[box] (Sigma) at (0,0) {$\Sigma$\\observable signals};
\node[box] (psi) at (4,0) {$\psi$\\encoding};
\node[box] (M) at (8,0) {$M$\\latent manifold};
\node[box] (Phi) at (8,2) {$\Phi$\\scalar estimate};
\node[box] (v) at (11,0.9) {$\mathbf{v}$\\flow estimate};
\node[box] (S) at (11,-0.9) {$S$\\entropy estimate};
\draw[arrow] (Sigma)--(psi);
\draw[arrow] (psi)--(M);
\draw[arrow] (M)--(Phi);
\draw[arrow] (M)--(v);
\draw[arrow] (M)--(S);
\end{tikzpicture}
\caption{Empirical observability schema. Observable signals are
mapped into a latent representation manifold, from which scalar,
vector, and entropy field estimates are constructed.}
\label{fig:observability-schema}
\end{figure}

\section{Estimating the RSVP Fields}

In neural data, $\Phi$ corresponds to activation magnitude; $\mathbf{v}$
is estimated from velocity vectors between successive population
states $x_{i+1} - x_i/(t_{i+1}-t_i)$; $S$ is estimated from the
entropy of the local state distribution. In linguistic systems, $\Phi$
is word frequency, $\mathbf{v}$ is the velocity of word embeddings
across diachronic corpora, and $S$ is contextual diversity.

\section{Identifying Amplitwist Transformations}

The local Jacobian $J_x \approx \alpha_x R_x$ can be estimated from
empirical trajectories by fitting the rotation and scaling components.
Neural jPCA analysis \cite{dicarlo2015} detects rotational dynamics
in population activity, providing a direct empirical handle on the
rotation component $R_x$.

\section{Limitations and the Realization Condition}

The realization condition (\cref{def:realization}) is the formal
bridge between theory and data: empirical methods estimate $\psi$ and
bound the tolerance $\varepsilon$. Measurement noise, partial
observability, and the requirement to approximate $\psi$ from finite
samples all contribute to the irreducible tolerance floor.

%%% CHAPTER 18 %%%

\chapter[Representation, Action, and Alignment]{Representation, Action, and the Problem of Alignment}
\label{ch:alignment}

\section{The Separation of Representation and Action}

A system possessing coherent representational structures does not
automatically produce actions that preserve those structures. This
chapter formalizes the representation--action gap categorically.

\section{Representation and Action Categories}

\begin{definition}[Representation Category]
\label{def:rep-category}
The representation category $\mathcal{M}$ has as objects the points
of $\mathcal{X}$, and as morphisms the entropy-modulated amplitwist
cascade operators $A^{(k)}$. Composition is cascade concatenation;
identity morphisms are $A = I$.
\end{definition}

\begin{definition}[Action Category]
The action category $\mathcal{A}$ has as objects action states and
morphisms transitions between actions.
\end{definition}

\begin{definition}[Action Functor]
A decision mechanism is a functor $G : \mathcal{M} \rightarrow
\mathcal{A}$ preserving composition and identities.
\end{definition}

\section{Alignment Degree and Generic Misalignment}

Let $\mathcal{M}_I \subseteq \mathcal{M}$ be the subcategory of
morphisms satisfying structural invariant $I$.

\begin{definition}[Alignment Degree]
$\alpha(G,I) = \mu(\{f \in \mathcal{M}_I : G(f)\text{ satisfies }
I\}) / \mu(\mathcal{M}_I)$.
\end{definition}

\begin{figure}[htbp]
\centering
\begin{tikzpicture}[>=Latex, every node/.style={font=\small}]
\draw[thick] (-4,-2) ellipse (2.3 and 1.3);
\node at (-4,2) {$\mathcal{M}$};
\draw[thick] (4,-2) ellipse (2.3 and 1.3);
\node at (4,2) {$\mathcal{A}$};
\draw[->,thick] (-2,-1)..controls(-0.5,-0.3)..(2,-1);
\draw[->,thick] (-2,0.6)..controls(-0.5,1.1)..(2,0.7);
\node at (0,1.6) {$G$};
\draw[dashed] (-4,-2) ellipse (0.9 and 0.6);
\node at (-4,-2.9) {$\mathcal{M}_I$};
\draw[dashed] (4,-2) ellipse (0.9 and 0.6);
\node at (4,-2.9) {$\mathcal{A}_I$};
\end{tikzpicture}
\caption{Alignment geometry. The functor $G$ maps representations to
actions. Alignment requires preservation of the invariant
subcategories $\mathcal{M}_I \to \mathcal{A}_I$.}
\label{fig:alignment-geometry}
\end{figure}

\begin{theorem}[Generic Misalignment]
\label{thm:generic-misalignment}
Let $\mathcal{F}(\mathcal{M},\mathcal{A})$ be the space of all functors
$G : \mathcal{M} \rightarrow \mathcal{A}$ with measure $\mu_F$
induced by the RSVP field structure. Then
\[
\mu_F(\{G : \alpha(G,I) = 1\}) = 0.
\]
\end{theorem}

\begin{proof}[Proof sketch]
$I$-structure-preserving functors satisfy the full-and-faithful
condition on $\mathcal{M}_I$, imposing a system of equalities on
the matrix entries of $G$. The set of such maps is a proper closed
submanifold of $\mathcal{F}(\mathcal{M},\mathcal{A})$, hence
measure zero with respect to any absolutely continuous measure $\mu_F$.
\end{proof}

\begin{figure}[htbp]
\centering
\begin{tikzpicture}[every node/.style={font=\small}]
\draw[thick] (0,0) ellipse (3.8 and 2.2);
\draw[thick] (1.2,0) ellipse (0.9 and 1.5);
\node at (0,2.55) {$\mathcal{F}(\mathcal{M},\mathcal{A})$};
\node at (1.2,0) {\shortstack{aligned\\functors}};
\node at (-1.7,-1.8) {generic functors};
\end{tikzpicture}
\caption{Measure-zero intuition for generic misalignment. The space
of aligned functors is a proper closed submanifold of the full functor
space.}
\label{fig:generic-misalignment-geometry}
\end{figure}

\section{Goodhart's Law as Functorial Failure}

\begin{proposition}
Let $\Phi_W, \Phi_P : M \rightarrow \mathbb{R}$. If $G$ maximizes
$\Phi_P$ rather than $\Phi_W$, then $G$ fails to preserve the
ordering invariant $I_W$ on any region where $\Phi_P$ and $\Phi_W$
conflict.
\end{proposition}

\section{Structure-Preserving Functors}

\begin{definition}[Structure-Preserving Functor]
$G$ is $I$-structure-preserving if $G|_{\mathcal{M}_I}$ is full
and faithful.
\end{definition}

By \cref{thm:generic-misalignment}, such functors require deliberate
design: they cannot arise by default from arbitrary training
procedures.

%%% CHAPTER 19 %%%

\chapter[Unified Constraint Geometry]{Unified Constraint Geometry: From Hierarchies to Aligned Intelligence}
\label{ch:synthesis}

\section{Overview}

This chapter synthesizes all preceding developments into a unified
mathematical statement and identifies the three principal results
of the monograph.

\begin{figure}[htbp]
\centering
\begin{tikzpicture}[>=Latex, every node/.style={font=\small, align=center},
  box/.style={draw, rounded corners, minimum width=2.7cm,
    minimum height=0.95cm}]
\node[box] (local) at (0,0) {local observations\\and constraints};
\node[box] (disc) at (4,1.6) {discrete structures\\trees, triplets};
\node[box] (geom) at (4,0) {geometric structures\\ultrametrics, cones};
\node[box] (field) at (4,-1.6) {field structures\\$(\Phi,\mathbf{v},S)$};
\node[box] (global) at (8,0) {global coherent\\structure};
\node[box] (action) at (12,0) {aligned action\\via preserved invariants};
\draw[->,thick] (local)--(disc);
\draw[->,thick] (local)--(geom);
\draw[->,thick] (local)--(field);
\draw[->,thick] (disc)--(global);
\draw[->,thick] (geom)--(global);
\draw[->,thick] (field)--(global);
\draw[->,thick] (global)--(action);
\end{tikzpicture}
\caption{Synthesis of the monograph. Local constraints generate
discrete, geometric, and field-theoretic structures that converge to
coherent global organization, which must then be preserved across the
representation--action boundary.}
\label{fig:unified-constraint-geometry}
\end{figure}

\section{Unified Constraint Reconstruction}

\begin{theorem}[Unified Constraint Reconstruction]
\label{thm:unified}
Let $\{C_i\}$ be a consistent and complete collection of constraints
on a state space $\mathcal{X}$. Then the entropy of any posterior
distribution over $\mathcal{X}$ consistent with $\{C_i\}$ converges
to zero monotonically in the ordering of constraints by information
gain.
\end{theorem}

\begin{proof}
By completeness (\cref{def:complete-constraint}), each constraint
eliminates at least one element not eliminated by previous constraints
until only the unique element $x^\ast$ remains. By
\cref{prop:entropy-decrease} each such step strictly reduces entropy.
The Triplet Reconstruction Theorem (\cref{thm:triplet-reconstruction}),
the Tree--Ultrametric Equivalence (\cref{thm:tree-ultrametric}), and
the constraint framework of \cref{ch:information} are special cases
corresponding to different choices of $\mathcal{X}$ and $\{C_i\}$.
\end{proof}

Structure emerges as the intersection of constraint manifolds. This is
the monograph's central thesis, now proved.

\section{Summary of Principal Results}

The Triplet Reconstruction Theorem (\cref{thm:triplet-reconstruction})
establishes that global hierarchical structure is determined by local
triplet constraints. The Cascade Stability Theorem
(\cref{thm:cascade-stability}) establishes that recursive amplitwist
transformations generate bounded convergent vorticity under the
explicitly characterized threshold $\epsilon_\mathrm{crit} = 1/(CN)$.
The Generic Misalignment Theorem (\cref{thm:generic-misalignment})
establishes that aligned functors form a measure-zero subset of all
decision maps.

Together these results instantiate \cref{thm:unified}: in each of the
discrete, continuous, and categorical settings treated here, structure
emerges through the coherent intersection of local constraints.

\section{Implications for AI and Future Directions}

Effective learning systems should emphasize mechanisms propagating
structural constraints through internal representations (by
\cref{thm:cascade-stability}). Stability requires entropy-like
regulatory mechanisms (by \cref{prop:entropy-modulation}).
Alignment requires deliberate $I$-structure-preserving design (by
\cref{thm:generic-misalignment}). Three open problems remain: a tight
characterization of $\epsilon_\mathrm{crit}$ via spectral properties
of $\{T_j\}$; a complete proof of \cref{thm:generic-misalignment} in
infinite-dimensional functor spaces; and empirical extraction of the
encoding map $\psi$ from neural population recordings.

\section{Conclusion}

Intelligence may be understood as the capacity of a system to navigate
constraint geometries, constructing coherent structures from incomplete
information while ensuring that those structures guide action in a
consistent and stable manner.

\backmatter

\chapter*{Afterword}

The theories presented here remain exploratory and incomplete. The
Cascade Stability Theorem provides a sufficient condition that is not
tight. The Generic Misalignment proof operates at the level of a
finite-dimensional linear map argument and requires extension to the
full infinite-dimensional functor space. The realization condition for
cortical columns requires methodological development for extracting
encoding maps from neural recordings. The stochastic RSVP framework
opens several analytic questions regarding the spectrum of the
Fokker--Planck operator and its dependence on the coupling parameters.
The sheaf-cohomological obstruction theory for constraint systems
admits a more systematic treatment via derived functors.

The categorical formalization of alignment failure as functorial
non-preservation of structural invariants may offer a more tractable
target for formal verification than behavioral or outcome-based
alignment criteria. Future investigations may extend the framework to
higher-dimensional and non-Euclidean manifolds, explore empirical
applications in neuroscience and machine learning, and refine the
coupling between representational and action-domain RSVP fields.

\appendix

\chapter{Notation and Symbol Reference}

\begin{center}
\renewcommand{\arraystretch}{1.2}
\begin{tabular}{lll}
\hline
Symbol & Meaning & First Introduced \\
\hline
$\mathcal{X}$ & State space & \cref{ch:foundations} \\
$X$ & Finite leaf set & \cref{ch:hierarchy} \\
$T^\ast$ & True latent tree & \cref{ch:hierarchy} \\
$\mathcal{R}(T)$ & Triplet system of $T$ & \cref{ch:hierarchy} \\
$\mathcal{H}$ & Hypothesis space & \cref{ch:hierarchy} \\
$C_o$ & Constraint operator & \cref{ch:information} \\
$d_T$ & Tree-induced ultrametric & \cref{ch:ultrametric} \\
$\mathcal{U}(X)$ & Space of ultrametrics on $X$ & \cref{ch:ultrametric} \\
$(X,d,\mu)$ & Metric measure space & \cref{ch:metric-measure} \\
$d_{GH}$ & Gromov--Hausdorff distance & \cref{ch:metric-measure} \\
$H(p)$ & Shannon entropy & \cref{ch:information} \\
$D_{KL}(p\|q)$ & KL divergence & \cref{ch:information} \\
$g_{ij}$ & Fisher information metric & \cref{ch:information} \\
$F(p)$ & Free energy functional & \cref{ch:information} \\
$M$ & Representation manifold & \cref{ch:rsvp} \\
$\Phi$ & Scalar field (semantic salience) & \cref{ch:rsvp} \\
$\mathbf{v}$ & Vector field (conceptual flow) & \cref{ch:rsvp} \\
$S$ & Entropy density field & \cref{ch:rsvp} \\
$\mathcal{E}$ & RSVP energy functional & \cref{ch:rsvp} \\
$G = \mathbb{R}_{>0}\times SO(n)$ & Amplitwist Lie group & \cref{ch:amplitwist-lie} \\
$\mathfrak{g}$ & Amplitwist Lie algebra & \cref{ch:amplitwist-lie} \\
$A_x$ & Amplitwist operator at $x$ & \cref{ch:amplitwist} \\
$A^{(\beta)}_x$ & Entropy-modulated amplitwist & \cref{ch:amplitwist} \\
$\varepsilon$ & Realization tolerance & \cref{ch:amplitwist} \\
$\mathcal{R}_k$ & $k$-layer deformation & \cref{ch:cascades} \\
$A^{(k)}$ & Layer-$k$ cascade amplitwist & \cref{ch:cascades} \\
$w_k$ & Entropy reliability weight & \cref{ch:cascades} \\
$\xi^{(N)}$ & Cascade vorticity & \cref{ch:cascades} \\
$\epsilon_\mathrm{crit}$ & Critical deformation bound & \cref{ch:amplitwist-lie} \\
$\mathcal{F}(U)$ & Sheaf sections over $U$ & \cref{ch:sheaf-constraints} \\
$L$ & Constraint Laplacian & \cref{ch:spectral-constraints} \\
$\gamma$ & Spectral gap & \cref{ch:spectral-constraints} \\
$P^\ast$ & Stationary distribution & \cref{ch:stochastic-rsvp} \\
$\mathcal{M}$ & Representation category & \cref{ch:alignment} \\
$\mathcal{A}$ & Action category & \cref{ch:alignment} \\
$G$ & Action functor & \cref{ch:alignment} \\
$\mathcal{M}_I$ & Invariant subcategory & \cref{ch:alignment} \\
$\alpha(G,I)$ & Alignment degree & \cref{ch:alignment} \\
\hline
\end{tabular}
\end{center}

\begin{thebibliography}{99}

\bibitem{needham1997}
T.~Needham.
\newblock \emph{Visual Complex Analysis}.
\newblock Oxford University Press, 1997.

\bibitem{mountcastle1957}
V.~B. Mountcastle.
\newblock Modality and topographic properties of single neurons of cat's
  somatic sensory cortex.
\newblock \emph{Journal of Neurophysiology}, 20(4):408--434, 1957.

\bibitem{horton2005}
J.~C. Horton and D.~L. Adams.
\newblock The cortical column: a structure without a function.
\newblock \emph{Philosophical Transactions of the Royal Society B},
  360(1456):837--862, 2005.

\bibitem{aho1981inferring}
A.~V. Aho, Y.~Sagiv, T.~G. Szymanski, and J.~D. Ullman.
\newblock Inferring a tree from lowest common ancestors.
\newblock \emph{SIAM Journal on Computing}, 10(3):405--421, 1981.

\bibitem{carhart2014}
R.~L. Carhart-Harris et al.
\newblock The entropic brain.
\newblock \emph{Frontiers in Human Neuroscience}, 8:20, 2014.

\bibitem{friston2010}
K.~Friston.
\newblock The free-energy principle: a unified brain theory?
\newblock \emph{Nature Reviews Neuroscience}, 11(2):127--138, 2010.

\bibitem{richman2000}
J.~S. Richman and J.~R. Moorman.
\newblock Physiological time-series analysis using approximate entropy and
  sample entropy.
\newblock \emph{American Journal of Physiology}, 278(6):H2039--H2049, 2000.

\bibitem{higgins2017}
I.~Higgins et al.
\newblock $\beta$-{VAE}: Learning basic visual concepts with a constrained
  variational framework.
\newblock \emph{ICLR}, 2017.

\bibitem{dicarlo2015}
J.~J. DiCarlo and D.~D. Cox.
\newblock Neural manifolds for the representation of learned motor skills.
\newblock \emph{Nature Neuroscience}, 18(4):512--520, 2015.

\bibitem{cohen2016}
T.~S. Cohen and M.~Welling.
\newblock Group equivariant convolutional networks.
\newblock \emph{ICML}, 2016.

\bibitem{connes1994}
A.~Connes.
\newblock \emph{Noncommutative Geometry}.
\newblock Academic Press, 1994.

\bibitem{vaswani2017}
A.~Vaswani et al.
\newblock Attention is all you need.
\newblock \emph{NeurIPS}, 30, 2017.

\bibitem{schwartz1980}
E.~L. Schwartz.
\newblock Computational anatomy and functional architecture of striate cortex.
\newblock \emph{Vision Research}, 20(8):645--669, 1980.

\bibitem{gu2004}
X.~Gu et al.
\newblock Genus zero surface conformal mapping and its application to brain
  surface mapping.
\newblock \emph{IEEE Transactions on Medical Imaging}, 23(8):949--958, 2004.

\bibitem{mackay2003}
D.~J.~C. MacKay.
\newblock \emph{Information Theory, Inference, and Learning Algorithms}.
\newblock Cambridge University Press, 2003.

\bibitem{amari2000}
S.~Amari and H.~Nagaoka.
\newblock \emph{Methods of Information Geometry}.
\newblock American Mathematical Society, 2000.

\bibitem{maclane1998}
S.~Mac Lane.
\newblock \emph{Categories for the Working Mathematician}.
\newblock 2nd~ed. Springer, 1998.

\bibitem{burago2001}
D.~Burago, Y.~Burago, and S.~Ivanov.
\newblock \emph{A Course in Metric Geometry}.
\newblock American Mathematical Society, 2001.

\bibitem{gromov1999}
M.~Gromov.
\newblock \emph{Metric Structures for Riemannian and Non-Riemannian Spaces}.
\newblock Birkh\"{a}user, 1999.

\bibitem{villani2009}
C.~Villani.
\newblock \emph{Optimal Transport: Old and New}.
\newblock Springer, 2009.

\bibitem{hall2015}
B.~Hall.
\newblock \emph{Lie Groups, Lie Algebras, and Representations}.
\newblock 2nd~ed. Springer, 2015.

\bibitem{bredon1997}
G.~E. Bredon.
\newblock \emph{Sheaf Theory}.
\newblock 2nd~ed. Springer, 1997.

\bibitem{chung1997}
F.~R.~K. Chung.
\newblock \emph{Spectral Graph Theory}.
\newblock American Mathematical Society, 1997.

\bibitem{risken1989}
H.~Risken.
\newblock \emph{The Fokker--Planck Equation}.
\newblock 2nd~ed. Springer, 1989.

\end{thebibliography}

\printindex

\end{document}
